\chapter{通信协议与基础控制模块设计}

\section{通信与基础控制设计目标}

总体架构确定了界面层、协调层、测试核心和串口通信层之间的职责边界。进入具体模块设计时，串口通信层需要完成从上层控制意图到硬件控制帧的转换，并将舵机控制器返回的反馈帧整理为测试流程能够使用的结构化样本。对于精密舵机测试软件而言，通信模块的功能不宜停留在字节收发层面，还应覆盖协议解释、帧完整性判断、连接状态维护和反馈数据组织。伺服测试平台与桌面测控软件的相关研究表明，执行机构测试过程通常同时包含设备控制、数据采集、状态显示和结果处理等环节，通信层边界会直接影响测试流程的可靠性\cite{li2014servotest,zhang2017qtflowmeter}。

基础控制模块位于串口通信层和测试核心之间。界面层给出的控制量先经协调层转换为协议帧，再通过 RS422 链路发送至舵机控制器；控制器返回的反馈帧则需要经过同步字识别、帧长判断、字段解析和校验验证，最终形成上层可复用的反馈样本。该链路贯穿串口连接、使能检测、手动控制、自动控制以及后续性能测试，是正式测试流程能够成立的基础。

根据上述约束，通信与基础控制模块的设计目标包括三点。第一，协议差异应限制在通信层内部，上层测试流程只面向统一的控制量和反馈量表达。第二，串口收发过程应具备顺序性和可追踪性，使控制帧发送、反馈帧接收和错误状态能够进入同一事件流。第三，基础控制过程需要区分手动发送和自动发送两类来源，避免多个任务同时占用同一串口发送通道。Electron 桌面数据采集应用的研究也说明，本地通信能力与界面逻辑保持清晰边界，有助于降低复杂桌面测控软件的维护成本\cite{gao2025electron}。

\section{RS422 串口链路设计}

所设计的软件系统通过 RS422 串口链路连接舵机控制器。串口管理模块承担端口枚举、连接建立、连接断开、数据写入、状态查询和异常恢复等职责。系统支持 115200 和 460800 两种波特率，数据位、停止位和校验方式分别固定为 8 位、1 位和无校验。该设置使界面层只需提供端口、波特率和协议版本等业务参数，底层串口参数则由通信模块统一配置。

串口连接过程按照端口枚举、配置选择、连接建立和状态广播的顺序执行。端口枚举为界面层提供可连接设备列表；连接建立后，系统保存当前串口配置，并向界面层发送连接成功事件；连接断开或串口错误发生时，协调层将错误信息转化为状态事件，供界面提示和任务中止使用。若断开并非由用户主动触发，系统根据已保存的串口配置尝试自动重连，默认重连间隔为 3000 ms。用户主动关闭串口时，系统取消重连计时，使界面操作优先于自动恢复过程。

发送前的连接状态检查是串口链路设计中的必要约束。手动控制、使能检测和自动控制均依赖同一物理发送通道，若串口未打开而继续生成控制帧，后续任务只能得到无意义的发送失败结果。系统因此在每次发送前检查连接状态，并在未连接时返回错误。该处理将链路可用性固定为测试任务启动前的必要条件，使通信异常能够在测试流程前端被发现。

\section{协议抽象与帧结构设计}

系统当前支持两种协议版本。协议一采用固定长度帧，协议二在控制字段之外增加帧长字段和序号字段。对于基于 RS422 的测试链路而言，通信设计不仅要描述帧结构，还要兼顾误码检测、数据采集和上位机联调等工程要求\cite{hu2026rs422ber}。在本系统中，两种协议共用同步字设计：计算机到控制器方向使用 0xAA55，对应字节序列 AA 55；控制器到计算机方向使用 0xAA5F，对应字节序列 AA 5F。

表~\ref{tab:protocol-frame-structure} 给出了两种协议的核心帧结构。协议一的控制帧和反馈帧均为 12 字节，适合固定长度的基础控制和反馈传输。协议二的控制帧为 18 字节，反馈帧为 26 字节，帧中包含帧长字段和序号字段，可在连续控制过程中标记帧次。两种协议的字段差异由通信层统一处理，上层测试流程只接收结构化后的控制量、反馈量和校验结果。

\begin{table}[htb]
  \centering
  \caption{协议帧结构对比}
  \label{tab:protocol-frame-structure}
  \begingroup
  \setlength{\tabcolsep}{3pt}
  \footnotesize
  \begin{tabular}{@{}p{0.10\textwidth}p{0.12\textwidth}p{0.10\textwidth}p{0.29\textwidth}p{0.13\textwidth}p{0.11\textwidth}@{}}
    \toprule
    协议 & 方向 & 帧长 & 字段顺序 & 校验区间 & 控制范围 \\
    \midrule
    协议一 & 计算机到控制器 & 12 字节 & 同步字、四路控制量、校验和 & 字节 0--9 & $\pm25^\circ$，即 $\pm2500$ \\
    协议一 & 控制器到计算机 & 12 字节 & 同步字、四路反馈量、校验和 & 字节 0--9 & 不适用 \\
    协议二 & 计算机到控制器 & 18 字节 & 同步字、帧长字段、序号、四路控制量、校验和 & 字节 2--15 & $\pm30^\circ$，即 $\pm3000$ \\
    协议二 & 控制器到计算机 & 26 字节 & 同步字、帧长字段、序号、四路控制量、四路反馈量、校验和 & 字节 2--23 & 不适用 \\
    \bottomrule
  \end{tabular}
  \endgroup
\end{table}

协议帧校验采用字节累加方式。设参与校验的字节为 $b_i$，起始字节和结束字节分别为 $s$ 和 $e$，则校验值为
\begin{equation}
  C=\left(\sum_{i=s}^{e} b_i\right)\bmod 2^{16}.
\end{equation}
其中，协议一对字节 0--9 累加，协议二控制帧对字节 2--15 累加，协议二反馈帧对字节 2--23 累加。串口通信研究表明，分包、缓存和校验机制对连续数据流的稳定传输具有直接影响\cite{wu2025serialpacket}。本系统将校验结果作为解析帧的一部分保存，上层模块既可使用反馈数据，也可判断该帧是否通过完整性检查。

解析后的帧被统一表示为帧头、载荷、校验、原始十六进制序列和时间戳。帧头保存协议版本、方向、同步字、帧长和序号；载荷保存四路控制量或四路反馈量；校验对象保存计算值、帧内期望值和验证结果。通过这种结构化表达，协议字段不再向测试算法外溢，时域测试、频域测试和数据显示模块均可基于统一反馈样本工作。

\section{串口发送与接收流程设计}

串口发送流程从控制值规范化开始。四路控制量以角度放大 100 倍后的整数形式进入通信层，协议一将控制范围限制在 $[-2500,2500]$，协议二将控制范围限制在 $[-3000,3000]$。控制值完成限幅后，通信层按照当前协议版本写入同步字、控制字段、帧长字段、序号字段和校验字段，并将完整帧加入发送队列。

发送队列用于保证串口写入顺序。手动控制、使能检测、自动控制和测试任务都可能产生控制帧，若多个来源直接写入串口，控制命令之间的顺序关系将难以确定。系统采用队列方式串行处理发送任务，并记录发送总数、失败总数、队列长度和最近发送时间。对于协议二，发送模块还维护递增序号，使连续控制帧具有可追踪的帧次标识。

接收流程面向连续字节流。串口接收到的数据先写入环形缓冲区，解析过程从缓冲区中搜索同步字。若找到同步字，系统再读取帧长字段判断协议类型和目标帧长；若数据不足一帧，解析过程暂停并等待后续字节；若缓冲区中连续包含多个完整帧，解析过程会连续提取。该方法可以同时处理半包和粘包问题，并减少无效字节对后续解析的影响。串口收发与帧解析流程如图~\ref{fig:serial-frame-flow} 所示。

\begin{figure}[htb]
  \centering
  \resizebox{\textwidth}{!}{%
  \begin{tikzpicture}[
    >=latex,
    block/.style={draw, align=center, minimum height=0.85cm, text width=2.3cm, font=\small},
    arrow/.style={->, thick}
  ]
    \node[block] (rx) at (0, 1.6) {串口接收\\连续字节流};
    \node[block] (buf) at (3.0, 1.6) {环形缓冲区\\追加数据};
    \node[block] (sync) at (6.0, 1.6) {同步字搜索\\定位帧头};
    \node[block] (len) at (9.0, 1.6) {帧长判断\\协议识别};
    \node[block] (extract) at (12.0, 1.6) {完整帧提取\\移出缓冲区};
    \node[block] (parse) at (12.0, 0) {字段解析\\控制量与反馈量};
    \node[block] (check) at (9.0, 0) {校验验证\\有效性标记};
    \node[block] (event) at (6.0, 0) {事件回传\\结构化样本};
    \node[block] (wait) at (3.0, 0) {等待后续字节\\保留残帧};

    \draw[arrow] (rx) -- (buf);
    \draw[arrow] (buf) -- (sync);
    \draw[arrow] (sync) -- (len);
    \draw[arrow] (len) -- node[above, font=\footnotesize]{完整} (extract);
    \draw[arrow] (extract) -- (parse);
    \draw[arrow] (parse) -- (check);
    \draw[arrow] (check) -- (event);
    \draw[arrow] (len) .. controls (8.4, -0.9) and (4.0, -0.9) .. node[below, font=\footnotesize]{数据不足} (wait);
    \draw[arrow] (wait) -- (buf);
  \end{tikzpicture}}
  \caption{串口收发与帧解析流程示意图}
  \label{fig:serial-frame-flow}
\end{figure}

完整帧提取后，解析模块根据协议版本和方向读取控制字段、反馈字段、序号字段和校验字段。若校验失败，帧仍可形成结构化结果，但校验状态被标记为无效；若帧格式无法识别，协调层向界面层发送解析错误事件。该处理方式区分了“帧可解析但校验失败”和“帧结构无法识别”两类情况，便于后续判断问题来自链路噪声、协议配置还是数据错位。

串口接收流程还承担反馈样本更新任务。控制器到计算机方向的有效反馈帧会被转换为四路反馈值，并保存为最近反馈样本。使能检测、预测试、时域测试和频域测试均可通过协调层读取该样本。由此，通信层既为界面提供实时显示数据，也为后续测试算法提供统一的数据入口。

\section{舵机使能检测算法设计}

使能检测用于在正式测试前判断四路执行对象是否对控制激励产生有效响应。若仅依据反馈绝对值设置阈值，静态偏置和传感器零漂可能造成误判。系统因此采用基线建模、小扰动激励和相关判别相结合的方法，在较小机械扰动条件下完成四路状态判定。

表~\ref{tab:enable-detection-params} 给出了使能检测的默认参数。检测过程采用 $2^\circ$、$4^\circ$、$6^\circ$ 三个幅值等级逐步试探，并限制单步控制量变化。该参数设置通过多次小扰动积累统计证据，降低单次大幅动作对机械状态和判定结果的影响。

\begin{table}[htb]
  \centering
  \caption{使能检测默认参数}
  \label{tab:enable-detection-params}
  \small
  \begin{tabular}{@{}p{0.26\textwidth}p{0.24\textwidth}p{0.40\textwidth}@{}}
    \toprule
    参数 & 默认值 & 作用 \\
    \midrule
    幅值等级 & $2^\circ$、$4^\circ$、$6^\circ$ & 逐级试探激励强度 \\
    控制间隔 & 10 ms & 保持采样和控制节奏一致 \\
    预热时间 & 180 ms & 发送零控制并稳定初始状态 \\
    基线时间 & 120 ms & 估计偏置和噪声 \\
    终判阈值 & $Z_{th}=4$ & 无早停时的显著性阈值 \\
    强通过阈值 & $Z_{hi}=6$ & 提前确认有效响应 \\
    强失败阈值 & $Z_{lo}=1.5$ & 提前确认无有效响应 \\
    增益阈值 & $g_{th}=0.15$ & 响应强度门槛 \\
    最大单步变化 & $0.8^\circ$ & 限斜率保护 \\
    \bottomrule
  \end{tabular}
\end{table}

检测开始后，系统先在零控制条件下采集基线样本。设第 $i$ 路反馈序列为 $f_i(m)$，基线窗口内样本数为 $M$，则偏置估计为
\begin{equation}
  b_i=\frac{1}{M}\sum_{m=1}^{M}f_i(m).
\end{equation}
后续响应分析使用去偏置反馈
\begin{equation}
  y_i(k)=f_i(k)-b_i.
\end{equation}
基线噪声标准差按下式估计：
\begin{equation}
  \sigma_i=\max\left(1,\sqrt{\frac{1}{M}\sum_{m=1}^{M}\left(f_i(m)-b_i\right)^2}\right).
\end{equation}
式中，下界 1 用于避免噪声估计过小时导致显著性指标异常放大。

基线建立后，系统对四路舵机施加分级编码激励。每级激励包含 16 组正负组合，该数量与基线采样规模保持同一量级，便于形成可用的相关统计量；四轴组合编码用于降低轴间响应耦合对判定结果的影响。为减少起停冲击，期望激励经过 $\sin^2$ 包络处理，并采用限斜率约束。若期望激励为 $u_i^\ast(k)$，实际输出控制量为 $u_i(k)$，则
\begin{equation}
  u_i(k)=u_i(k-1)+\operatorname{clip}\left(u_i^\ast(k)-u_i(k-1),-\Delta_{\max},\Delta_{\max}\right),
\end{equation}
其中，$\Delta_{\max}$ 对应 $0.8^\circ$ 的最大单步变化。

在长度为 $K$ 的窗口内，第 $i$ 路激励与反馈之间的相关项和激励能量定义为
\begin{equation}
  R_i(K)=\sum_{k=1}^{K}u_i(k)y_i(k),\qquad E_i(K)=\sum_{k=1}^{K}u_i^2(k).
\end{equation}
由此得到增益估计和噪声归一化显著性指标：
\begin{equation}
  \hat g_i(K)=\frac{R_i(K)}{E_i(K)},\qquad
  Z_i(K)=\frac{|R_i(K)|}{\max\left(1,\sigma_i\sqrt{\max(1,E_i(K))}\right)}.
\end{equation}
其中，$R_i$ 表示反馈与激励的同调程度，$\hat g_i$ 表示响应强度，$Z_i$ 表示经过噪声归一化后的显著性。若 $Z_i\ge Z_{hi}$ 且 $|\hat g_i|\ge g_{th}$，该轴可提前判定为有效；若采样数满足要求且 $Z_i\le Z_{lo}$，该轴可提前判定为无有效响应。未触发早停时，终判规则为
\begin{equation}
  s_i=
  \begin{cases}
    \mathrm{ENABLED}, & Z_i\ge Z_{th}\ \land\ |\hat g_i|\ge g_{th},\\
    \mathrm{DISABLED}, & \text{otherwise}.
  \end{cases}
\end{equation}

使能检测算法流程如图~\ref{fig:enable-detection-flow} 所示。检测结束后，系统按照限斜率约束将控制量逐步回到零值，并清空发送队列。主进程侧最多执行 3 次检测尝试；若连续尝试仍无法获得有效结果，则返回检测失败。检测完成后控制模式恢复为手动模式，以免自动激励状态影响后续通信性能测试和正式性能测试。

\begin{figure}[htb]
  \centering
  \resizebox{0.98\textwidth}{!}{%
  \begin{tikzpicture}[
    >=latex,
    block/.style={draw, align=center, minimum height=0.85cm, text width=2.35cm, font=\small},
    arrow/.style={->, thick}
  ]
    \node[block] (warm) at (0, 1.8) {预热阶段\\零控制稳定};
    \node[block] (base) at (3.0, 1.8) {基线采样\\估计偏置与噪声};
    \node[block] (excite) at (6.0, 1.8) {分级激励\\组合编码输入};
    \node[block] (stat) at (9.0, 1.8) {相关统计\\计算增益与显著性};
    \node[block] (early) at (12.0, 1.8) {早停判断\\强通过或强失败};
    \node[block] (final) at (12.0, 0) {终判规则\\四路状态输出};
    \node[block] (zero) at (9.0, 0) {回零处理\\限斜率恢复};
    \node[block] (manual) at (6.0, 0) {恢复手动模式\\清理发送状态};

    \draw[arrow] (warm) -- (base);
    \draw[arrow] (base) -- (excite);
    \draw[arrow] (excite) -- (stat);
    \draw[arrow] (stat) -- (early);
    \draw[arrow] (early) -- node[right, font=\footnotesize]{满足} (final);
    \draw[arrow] (final) -- (zero);
    \draw[arrow] (zero) -- (manual);
    \draw[arrow] (early) .. controls (12.2, 3.3) and (6.0, 3.3) .. node[above, font=\footnotesize]{继续采样} (excite);
  \end{tikzpicture}}
  \caption{舵机使能检测算法流程示意图}
  \label{fig:enable-detection-flow}
\end{figure}

\section{基础控制与信号发生器设计}

基础控制模块负责协调手动发送、自动发送和周期性信号生成。手动控制由界面侧输入四路控制量，经过通信层构造协议帧后写入串口；自动控制由后台任务按照固定周期生成控制量，再通过同一发送队列进入串口。两类控制来源共享发送通道，因此系统设置手动和自动两种控制模式，并在自动模式下拒绝手动发送请求。

周期性信号生成以正弦波为基础。设幅值为 $A$、频率为 $f$、相位为 $\varphi$、偏置为 $d$，则时刻 $t$ 的目标角度为
\begin{equation}
  r(t)=d+A\sin(2\pi ft+\varphi).
\end{equation}
后台任务按照固定间隔计算 $r(t)$，并将角度转换为协议使用的整数控制量。当前自动控制周期默认取 10 ms。由于控制序列由后台统一调度，界面事件只表达启动、停止和参数配置意图，控制输出不依赖界面刷新节奏。

信号发生器模块在设计上分为波形采样、轨迹整形、状态机、调度器和适配器几部分。波形采样给出理想目标角度；轨迹整形用于限制速度和加速度变化；状态机约束 idle、running、checking、stopping 和 fault 等状态转换；调度器提供固定周期推进；适配器负责连接串口发送和反馈输入。这种分层使波形、状态和串口边界可以分别验证，并为后续频域扫频和更复杂的控制策略保留接口。

在本论文范围内，基础控制模块主要承担两项任务：其一，为界面手动控制和自动正弦控制提供统一发送链路；其二，为使能检测和后续测试任务提供可复用的控制输出机制。闭环控制律和更复杂的轨迹约束属于可扩展能力，不作为当前已经验证的主要功能展开。

\section{本章小结}

本章围绕 RS422 串口链路、协议抽象、串口收发、使能检测和基础控制建立了软件的底层通信与控制基础。系统通过协议一和协议二统一控制量与反馈量表达，通过发送队列、环形缓冲区和事件回传机制组织连续串口数据，再通过基线建模、小扰动激励和相关判别完成舵机使能检测。

上述设计把总体架构中的通信层具体化为可执行的协议帧、发送流程、接收流程和状态判定方法。后续测试功能可以在该基础上获得控制发送能力、反馈样本和连接状态，无需直接处理原始字节流或硬件初始化细节。由通信与基础控制模块提供的稳定数据入口，也为通信性能测试、控制性能预测试、时域性能测试和频域性能测试的算法设计提供了执行条件。
